% Generated draft: requires mathematical and source review.
\documentclass[11pt]{article}
\usepackage[margin=25mm]{geometry}
\usepackage{fontspec}
\setmainfont{Latin Modern Roman}
\usepackage{amsmath,amssymb,mathtools,unicode-math}
\setmathfont{Latin Modern Math}
\usepackage{enumitem,needspace,xurl,hyperref}
\setlength{\parindent}{0pt}
\setlength{\parskip}{5pt}
\setlength{\emergencystretch}{3em}
\newcommand{\N}{\mathbb{N}}
\newcommand{\Z}{\mathbb{Z}}
\newcommand{\Q}{\mathbb{Q}}
\newcommand{\R}{\mathbb{R}}
\newcommand{\C}{\mathbb{C}}
\newcommand{\F}{\mathbb{F}}
\newcommand{\A}{\mathbb{A}}
\newcommand{\PP}{\mathbb{P}}
\newcommand{\E}{\mathbb{E}}
\newcommand{\eps}{\varepsilon}
\newcommand{\dd}{\,\mathrm{d}}
\begin{document}
\section*{UnsolvedMath ID: 700792}
\textbf{Source problem code:} catalogue entry
\subsection{Definitions and mathematical statement}
A group \(G\) is \emph{periodic} if every element has finite order:
\[
\forall g\in G\;\exists n\in\Z_{\geq 1}\quad g^n=1.
\]
For a prime \(p\), a \(p\)-group here means a group in which every element has order a power of \(p\); finiteness of the group is not assumed. In particular, a \(2\)-group has this property for \(p=2\). An extension of a group \(N\) by a group \(Q\) is a group \(E\) occurring in a short exact sequence
\[
1\longrightarrow N\longrightarrow E\longrightarrow Q\longrightarrow 1.
\]

The supplied problem uses the terms \emph{conjugacy biprimitively finite group} and \emph{binary finite group} without defining them. It refers to Kourovka Notebook Problem 6.59 for the former terminology, but that referenced text is not supplied. Accordingly, let \(\mathcal C\) denote the class intended by the first term, let \(\mathcal B\) denote the class intended by the second term, and let \(\mathcal G_2\) denote the class of all \(2\)-groups. These symbols preserve unresolved source terminology; they do not supply its missing definitions.

\textbf{Question (S. P. Strunkov, 1990).} Is every periodic group \(G\in\mathcal C\) obtainable from groups in \(\mathcal G_2\cup\mathcal B\) by taking extensions? The source further singles out the same question restricted to \(p\)-groups, for each odd prime \(p\), as being of independent interest.

\textbf{Unresolved qualifications.} A complete formal statement requires the intended definitions of both named classes and the precise meaning of ``obtained \dots\ by taking extensions.'' In particular, the supplied text does not specify whether only finite iteration of extensions is intended or whether some additional infinite construction is permitted. No such construction convention is imposed in this draft.
\subsection{Short English statement}
Strunkov asks whether every periodic conjugacy biprimitively finite group can be constructed from \(2\)-groups and binary finite groups by taking extensions, highlighting the restriction to \(p\)-groups for odd primes \(p\). The supplied record omits essential terminology and construction conventions, so the formulation remains qualified.
\subsection{Sources}
\begin{itemize}
\item Local UnsolvedMath problem registry, permanent ID \(700792\), exact source problem code \texttt{catalogue entry}, title \emph{Kourovka Notebook Problem 11.100}; supplied source record, upstream revision \texttt{d371ee917ac9647fbb06cd6fc529deec6f3dda82}.
\item The supplied attribution is to S. P. Strunkov, proposed in 1990, in \emph{The Kourovka Notebook (Unsolved Problems in Group Theory)}, 11th issue (1990), Problem 11.100. The background also identifies the 21st issue edition (2026), \texttt{arXiv:1401.0300v48}. These bibliographic details are supplied attribution and have not been independently verified.
\item The statement refers to Problem 6.59 for terminology; the text of that problem is absent from the supplied record.
\item The registry labels the problem \texttt{open}. This is a supplied status label, not an independently verified assessment of the literature. The record identifies the collection as Issues 1--21 (1965--2026), edited by E. I. Khukhro and V. D. Mazurov.
\end{itemize}
\end{document}
